% ack.tex
% mainfile: perfbook.tex
% SPDX-License-Identifier: CC-BY-SA-3.0

\chapter{致谢}
\label{bck:ack:Credits}
%
\Epigraph{如果我看得更远，那是因为我站在巨人的肩膀上。}
	 {Isaac Newton，现代化表述}

% \section{Authors}

% CCASA 3.0 US wording from John Wiegley.
%  http://www.newartisans.com/blog_assets/git.from.bottom.up.pdf

\section{\LaTeX\ 顾问}

Akira Yokosawa 是本书的 \LaTeX\ 顾问，其工作最为突出的部分包括
维护\cref{chp:app:styleguide:Style Guide}中列出的风格指南。
这项工作包括表格布局、代码清单、字体、数学公式渲染、
缩略语、参考文献格式、题词、超链接和纸张大小。
Akira 还完善了快速测验的交叉引用，使得在快速测验和其答案之间
能够轻松准确地导航。
他还添加了构建选项，允许隐藏快速测验以及将其收集到每章末尾
（教科书风格）。

这一角色还包括构建系统，Akira 对其进行了优化并使其更加
用户友好。
他的改进包括自动响应参考文献变更、自动确定哪些源文件存在，
以及从源文件自动生成代码清单（带有自动生成的超链接行号引用）。

\section{审阅者}

\begin{itemize}
\item	Alan Stern（\cref{chp:Advanced Synchronization: Memory Ordering}）。
\item	Andy Whitcroft（\cref{sec:defer:RCU Fundamentals}，
	\cref{sec:defer:RCU Linux-Kernel API}）。
\item	Artem Bityutskiy（\cref{chp:Advanced Synchronization: Memory Ordering}，
	\cref{chp:app:whymb:Why Memory Barriers?}）。
\item	Dave Keck（\cref{chp:app:whymb:Why Memory Barriers?}）。
\item	David S. Horner
	（\cref{sec:formal:Promela Parable: dynticks and Preemptible RCU}）。
\item	Gautham Shenoy（\cref{sec:defer:RCU Fundamentals}，
	\cref{sec:defer:RCU Linux-Kernel API}）。
\item	``jarkao2''，即 LWN 访客 \#41960（\cref{sec:defer:RCU Linux-Kernel API}）。
\item	Jonathan Walpole（\cref{sec:defer:RCU Linux-Kernel API}）。
\item	Josh Triplett（\cref{chp:Formal Verification}）。
\item	Michael Factor（\cref{sec:future:Transactional Memory}）。
\item	Mike Fulton（\cref{sec:defer:RCU Fundamentals}）。
\item	Peter Zijlstra
	（\cref{sec:defer:RCU Usage}）。% Lanin and Shasha citation.
\item	Richard Woodruff（\cref{chp:app:whymb:Why Memory Barriers?}）。
\item	Suparna Bhattacharya（\cref{chp:Formal Verification}）。
\item	Vara Prasad
	（\cref{sec:formal:Promela Parable: dynticks and Preemptible RCU}）。
\end{itemize}

以补丁形式提供反馈的审阅者（这种形式极受欢迎）在 git 日志中致谢。

\section{机器提供者}

读者可能注意到一些展示了多达数百个 CPU 的可扩展性数据的图表，
这要感谢我目前的雇主，特别感谢 Paul Saab、Yashar Bayani、
Joe Boyd 和 Kyle McMartin。

回顾我在 IBM 的时光，
非常感谢 Martin Bligh，他创建了 IBM Linux 技术中心的
Advanced Build and Test (ABAT) 系统，
以及 Andy Whitcroft、Dustin Kirkland 和许多扩展该系统的人员。
同时也非常感谢众多机器提供者：
Andrew Theurer、
Andy Whitcroft、
Anton Blanchard、
Chris McDermott、
Cody Schaefer、
Darrick Wong、
David ``Shaggy'' Kleikamp、
Jon M. Tollefson、
Jose R. Santos、
Marvin Heffler、
Nathan Lynch、
Nishanth Aravamudan、
Tim Pepper
和
Tony Breeds。

\section{原始出版物}

\ListOriginalPublications

\section{图片致谢}

\ListContributions

\Cref{fig:defer:RCU Areas of Applicability}改编自
\ppl{Fedor}{Pikus} 的``When to use RCU''幻灯片~\cite{FedorPikus2017RCUthenWhat}。
\cref{sec:defer:Reference Counting}中关于机械引用计数器的讨论
源于与 \ppl{Dave}{Regan} 的私人交谈。

\section{其他支持}

我们感谢许多 CPU 架构师耐心地解释其 CPU 的指令重排序和
内存重排序特性，特别是
Wayne Cardoza、Ed Silha、Anton Blanchard、Tim Slegel、Juergen Probst、
Ingo Adlung、Ravi Arimilli、Cathy May、Derek Williams、
H.~Peter Anvin、
Andy Glew、Leonid Yegoshin、
Richard Grisenthwaite 和 Will Deacon。
Wayne 特别值得感谢，他耐心地解释了 Alpha 对依赖加载的重排序——
Paul 曾非常顽固地抵制这个教训！

Association for Computing Machinery 的 bibtex 生成服务
为我们编写参考文献节省了大量时间和精力，对此我们深表感激。
同时也感谢 Stamatis Karnouskos，他说服我将我那古老的参考文献数据库
踢踢打打地拖入了21世纪。
这类技术工作都要感谢维持互联网和万维网运行的众多个人和组织，
本书也不例外。

本材料的部分内容基于美国国家科学基金会 Grant No.\@ CNS-0719851
资助的工作。
